\section{Results}
\label{sec:results}

Table~\ref{tab:main} is the study's core measurement: every seed-0 arm against
the seed-matched baseline, on the held-out selection metric, the external C4
set, and the dense paired training statistic, with its step-time and memory
cost. Figure~\ref{fig:plane} plots the same arms in the plane that matters for
deciding whether any of this is worth doing.

\begin{table}[t]
\centering
\caption{Seed-0 arms against the seed-matched input-side baseline, 0.5M-token
batch, full 2.68B-token budget. Deltas are means over the 11 held-out
evaluations in the second half of the budget; negative is better. $\Delta$ train
is the mean per-update paired difference over the final 536{,}870{,}912 tokens,
with its autocorrelation-inflated standard error. Step overhead is the median
over updates 2049--4096. Rows are ordered by selection delta. The two
\textbf{null} rows execute the candidate path and apply identity factors: they
are the reference any geometry claim must beat, not the baseline.}
\label{tab:main}
\footnotesize\setlength{\tabcolsep}{3pt}
% Generated by scripts/analyze.py. Do not edit.
\begin{tabular}{lrrrrrrr}
\toprule
Arm & $\Delta$ sel. & $\Delta$ C4 & $\Delta$ train & s.e. & Step \% & E2E \% & Peak GiB \\
\midrule
mlpdiag (leading) & \textbf{$-0.0066$} & $-0.0061$ & $-0.0051$ & $0.00004$ & $+21.9$ & $+19.3$ & 22.6 \\
dense, lagged & $-0.0059$ & $-0.0059$ & $-0.0057$ & $0.00004$ & $+25.6$ & $+31.3$ & 24.5 \\
dense, both widths & $-0.0058$ & $-0.0064$ & $-0.0050$ & $0.00004$ & $+39.2$ & $+44.2$ & 24.2 \\
mlpdiag, lagged & $-0.0057$ & $-0.0050$ & $-0.0050$ & $0.00004$ & $+18.0$ & $+16.7$ & 22.7 \\
all diagonal & $-0.0045$ & $-0.0036$ & $-0.0037$ & $0.00004$ & $\mathbf{+10.4}$ & $+8.2$ & 22.4 \\
outer factor only & $-0.0040$ & $-0.0048$ & $-0.0039$ & $0.00004$ & $+29.6$ & $+35.0$ & 24.2 \\
no 2048 output & $-0.0031$ & $-0.0026$ & $-0.0021$ & $0.00004$ & $+18.0$ & $+16.8$ & 21.5 \\
down only & $-0.0014$ & $-0.0008$ & $-0.0012$ & $0.00004$ & $+4.4$ & $+5.0$ & 20.6 \\
\addlinespace
\textbf{null (dense)} & $-0.0013$ & $-0.0023$ & $-0.0014$ & $0.00004$ & $+37.8$ & $+45.0$ & 24.2 \\
\textbf{null (mlpdiag)} & $-0.0006$ & $-0.0020$ & $-0.0007$ & $0.00004$ & $+21.6$ & $+20.4$ & 22.6 \\
\addlinespace
\textbf{MuonT} (base) & $0$ & $0$ & $0$ & --- & $0$ & $0$ & 20.5 \\
down, output replaces & $+0.0025$ & $+0.0036$ & $+0.0032$ & $0.00005$ & $+2.7$ & $+2.7$ & 19.2 \\
Muon & $+0.0055$ & $+0.0053$ & $+0.0062$ & $0.00016$ & $-1.9$ & $-1.6$ & 19.0 \\
output only & $+0.0087$ & $+0.0109$ & $+0.0079$ & $0.00012$ & $+33.6$ & $+41.0$ & 22.2 \\
both sides diagonal & $+0.0180$ & $+0.0177$ & $+0.0134$ & $0.00031$ & $+16.4$ & $+13.1$ & 21.4 \\
\bottomrule
\end{tabular}

\end{table}

\begin{figure}[t]
\centering
\begin{tikzpicture}
\begin{axis}[
  width=0.88\textwidth, height=8.4cm,
  xlabel={Median step-time overhead versus input-side MuonT (\%)},
  ylabel={$\Delta$ held-out selection loss (nats)},
  xmin=-6, xmax=52, ymin=-0.0088, ymax=0.021,
  grid=major, grid style={gray!18},
  axis line style={gray!55}, tick align=outside,
  every axis plot/.append style={mark size=2.1pt},
  scaled y ticks=false,
  yticklabel style={/pgf/number format/fixed, /pgf/number format/precision=3},
  legend columns=3,
  legend style={at={(0.5,-0.17)}, anchor=north,
                font=\footnotesize, draw=gray!40, fill=white, column sep=8pt},
  label style={font=\small}, tick label style={font=\footnotesize},
]
% Generated by scripts/render_results_viz.py. Do not edit.
% Shaded net-of-null improvement zone (below Null-Dense = -0.001286)
\fill[famWT!6] (axis cs:-5.5,-0.0082) rectangle (axis cs:44.5,-0.001286);
\draw[pagemuted!70, dashed, line width=0.55pt] (axis cs:-5.5,0) -- (axis cs:44.5,0);
\draw[pagemuted!45, dotted, line width=0.5pt] (axis cs:0,-0.0082) -- (axis cs:0,0.0108);
\draw[famNull!85, densely dashed, line width=0.7pt] (axis cs:-5.5,-0.001286) -- (axis cs:44.5,-0.001286);
\node[font=\tiny, text=famNull!90!black, anchor=south east] at (axis cs:44.0,-0.001286) {null floor ($-0.0013$)};
\addplot[only marks, mark=*, mark size=2.3pt, color=famWT] coordinates {(21.949,-0.006584) (17.952,-0.005683)};
\addlegendentry{Width-thresholded (\texttt{WT})}
\addplot[only marks, mark=diamond*, mark size=2.6pt, color=famOD] coordinates {(10.399,-0.004503)};
\addlegendentry{Pure-diagonal (\texttt{OD})}
\addplot[only marks, mark=pentagon*, mark size=2.4pt, color=famDT] coordinates {(39.164,-0.005782) (25.592,-0.005893)};
\addlegendentry{Dense two-sided (\texttt{DT})}
\addplot[only marks, mark=square*, mark size=2.1pt, color=famNull] coordinates {(21.628,-0.000649) (37.837,-0.001286)};
\addlegendentry{Matched nulls (\texttt{Null})}
\addplot[only marks, mark=triangle*, mark size=2.3pt, color=famAbl] coordinates {(4.385,-0.001370) (18.042,-0.003068) (29.605,-0.003977) (2.738,0.002512) (33.556,0.008707)};
\addlegendentry{Ablations (\texttt{Abl})}
\addplot[only marks, mark=oplus*, mark size=2.3pt, color=famRef] coordinates {(-1.933,0.005505) (0.000,0.000000)};
\addlegendentry{Baselines (\texttt{Ref})}
\addplot[only marks, mark=star, mark size=3.5pt, line width=0.9pt, color=famWT] coordinates {(1.810,-0.005683) (3.090,-0.005893)};
\addlegendentry{\texttt{fast\_lag} frontier}
\draw[famWT, line width=0.9pt] (axis cs:1.630,-0.005683) -- (axis cs:2.090,-0.005683);
\draw[famDT, line width=0.9pt] (axis cs:2.820,-0.005893) -- (axis cs:3.370,-0.005893);
\draw[-Latex, famWT, densely dashed, line width=0.85pt] (axis cs:17.402,-0.005683) -- (axis cs:2.460,-0.005683);
\draw[-Latex, famDT, densely dotted, line width=0.8pt] (axis cs:25.042,-0.005893) to[bend left=10] (axis cs:3.790,-0.005893);
\node[font=\tiny, text=famWT, anchor=north, yshift=-2.5pt, inner sep=1pt] at (axis cs:21.949,-0.006584) {\textbf{\taxid{WT-MLP}}};
\node[font=\tiny, text=famWT, anchor=north east, xshift=-1pt, yshift=-2pt, inner sep=1pt] at (axis cs:17.952,-0.005683) {\taxid{WT-MLP-Lag}};
\node[font=\tiny, text=famOD, anchor=west, xshift=3pt, inner sep=1pt] at (axis cs:10.399,-0.004503) {\textbf{\taxid{OD-All}}};
\node[font=\tiny, text=famDT!90!black, anchor=north, yshift=-2.5pt, inner sep=1pt] at (axis cs:39.164,-0.005782) {\taxid{DT-All}};
\node[font=\tiny, text=famDT!90!black, anchor=south west, xshift=1.5pt, yshift=1.5pt, inner sep=1pt] at (axis cs:25.592,-0.005893) {\taxid{DT-All-Lag}};
\node[font=\tiny, text=famNull!90!black, anchor=south, yshift=2.5pt, inner sep=1pt] at (axis cs:21.628,-0.000649) {\taxid{Null-WT}};
\node[font=\tiny, text=famNull!90!black, anchor=south, yshift=2.5pt, inner sep=1pt] at (axis cs:37.837,-0.001286) {\taxid{Null-Dense}};
\node[font=\tiny, text=famAbl!90!black, anchor=south west, xshift=2pt, yshift=1.5pt, inner sep=1pt] at (axis cs:4.385,-0.001370) {\taxid{Abl-DownOnly}};
\node[font=\tiny, text=famAbl!90!black, anchor=south west, xshift=2pt, yshift=1pt, inner sep=1pt] at (axis cs:18.042,-0.003068) {\taxid{Abl-No2048}};
\node[font=\tiny, text=famAbl!90!black, anchor=west, xshift=3pt, inner sep=1pt] at (axis cs:29.605,-0.003977) {\taxid{Abl-OuterDense}};
\node[font=\tiny, text=famAbl!90!black, anchor=west, xshift=3pt, inner sep=1pt] at (axis cs:2.738,0.002512) {\taxid{Abl-DownSwap}};
\node[font=\tiny, text=famAbl!90!black, anchor=east, xshift=-3pt, inner sep=1pt] at (axis cs:33.556,0.008707) {\taxid{Abl-OutOnly}};
\node[font=\tiny, text=famRef, anchor=west, xshift=3pt, inner sep=1pt] at (axis cs:-1.933,0.005505) {\taxid{Muon}};
\node[font=\tiny, text=famRef, anchor=south east, xshift=2pt, yshift=2pt, inner sep=1pt] at (axis cs:0.000,0.000000) {\taxid{MuonT}};
\draw[famWT!75, line width=0.45pt] (axis cs:1.810,-0.005463) -- (axis cs:0.2,-0.00325);
\node[font=\tiny, text=famWT, anchor=west, fill=white, fill opacity=0.88, text opacity=1, inner sep=1pt] at (axis cs:-5.1,-0.00295) {\textbf{\taxid{WT-MLP-Lag}} (\texttt{fast} $+1.8\%$)};
\draw[famDT!80, line width=0.45pt] (axis cs:3.090,-0.006113) -- (axis cs:1.8,-0.00730);
\node[font=\tiny, text=famDT!90!black, anchor=west, inner sep=1pt] at (axis cs:1.5,-0.00755) {\textbf{\taxid{DT-All-Lag}} (\texttt{fast} $+3.1\%$)};
\draw[-Latex, famAbl, line width=0.85pt] (axis cs:16.445,0.0093) -- (axis cs:16.445,0.01055);
\node[font=\tiny, text=famAbl!90!black, anchor=west, inner sep=1.5pt] at (axis cs:17.045,0.00995) {\taxid{Abl-BothDiag} ($+0.0180$)};

\draw[gray!55, dashed] (axis cs:-6,0) -- (axis cs:50,0);
\draw[laOrangeLine!70, dashed] (axis cs:-6,-0.001286) -- (axis cs:50,-0.001286);
\node[font=\tiny, text=laOrangeLine!80!black, anchor=west]
  at (axis cs:-5.5,-0.00205) {null level};
\end{axis}
\end{tikzpicture}
\caption{Cost and quality for the seed-0 arms. Down and to the left is better.
The dashed orange line is what a matched null achieves with \emph{identity}
factors, i.e.\ with no output geometry at all: it is the real zero for this
plot, and four of the eleven candidates do not clearly clear it. The cheapest
candidate that does, \alldiag{} at $+10.4$\% step time, retains about
two-thirds of the leading arm's gain for less than half its overhead. Nothing
appears in the lower-left quadrant, because no arm on this path is cheaper than
the baseline.}
\label{fig:plane}
\end{figure}

\subsection{The output side complements input whitening and cannot replace it}
\label{sec:complement}

The comparison that carries this claim is between two arms applying
\emph{identical} output factors to all 224 matrices, differing only in whether
the input factors stay. \dense{} keeps them and improves selection loss by
$0.005782$ nats; \leftall{} removes them and is $0.008707$ nats \emph{worse}
than the baseline. The swing between the two is 0.0145 nats, far larger than
either arm's distance from the baseline.

\begin{finding}[The two sides are not interchangeable]
The same output-side factors help by $0.0058$ nats when added to input
whitening and hurt by $0.0087$ nats when substituted for it. The arm that
drops input factors loses more than the arm that adds output factors gains, so
the output moment is not a cheaper restatement of the input moment's
information. On the narrow down-projection-only arms the same pattern holds:
adding an outer output factor gives $-0.001370$, while replacing the
down-projection input factors with output ones gives $+0.002512$.
\end{finding}

This asymmetry has a plausible structural reading that our data supports but
does not prove. The input moment conditions the \emph{domain} of the update,
which is what the orthogonalization in \eqref{eq:ns} then acts on; the output
moment rescales the \emph{codomain} after the fact. Removing the first leaves
$\Nten$ operating on an unconditioned matrix, and no amount of output-side
rescaling recovers that. Consistent with this, the fully diagonal arm
\texttt{inner\_outer\_all\_diag} --- the only arm that degrades the \emph{input}
side as well --- is the worst arm in the study at $+0.018011$, worse than plain
Muon, despite applying output factors in the same places as the leading arm.
Whatever value the output side has here, it is conditional on a well-conditioned
input side, and cheapening the input side is a much more expensive mistake than
cheapening the output side.

Figure~\ref{fig:traindelta} shows these orderings are not artifacts of the
evaluation window. The paired training-loss differences separate within the
first 200M tokens and keep their sign for the remaining 2.5B.

\begin{figure}[t]
\centering
\begin{tikzpicture}
\begin{axis}[
  width=0.88\textwidth, height=7.4cm,
  xlabel={Tokens processed (billions)}, ylabel={Paired $\Delta$ training loss (nats)},
  xmin=0, xmax=2.69, ymin=-0.020, ymax=0.016,
  grid=major, grid style={gray!18},
  axis line style={gray!55}, tick align=outside,
  scaled y ticks=false,
  yticklabel style={/pgf/number format/fixed, /pgf/number format/precision=3},
  legend columns=2,
  legend style={at={(0.5,-0.19)}, anchor=north,
                font=\scriptsize, draw=gray!40, fill=white, column sep=6pt,
                row sep=1pt},
  legend cell align=left,
  label style={font=\small}, tick label style={font=\footnotesize},
  every axis plot/.append style={line width=0.9pt},
]
\addplot[color=laBlueLine] table[x=tokens, y=inner_outer_mlpdiag] {figs/train_delta.dat};
\addlegendentry{mlpdiag (leading)}
\addplot[color=laBlueLine!55, densely dashed] table[x=tokens, y=inner_outer_all] {figs/train_delta.dat};
\addlegendentry{dense, both widths}
\addplot[color=laPurpleLine] table[x=tokens, y=inner_outer_leftdiag] {figs/train_delta.dat};
\addlegendentry{all diagonal}
\addplot[color=laOrangeLine] table[x=tokens, y=null_all] {figs/train_delta.dat};
\addlegendentry{null (dense)}
\addplot[color=laOrangeLine!60, densely dotted] table[x=tokens, y=null_mlpdiag] {figs/train_delta.dat};
\addlegendentry{null (mlpdiag)}
\addplot[color=crossred] table[x=tokens, y=left_all] {figs/train_delta.dat};
\addlegendentry{output only}
\addplot[color=crossred!60, densely dashed] table[x=tokens, y=inner_outer_all_diag] {figs/train_delta.dat};
\addlegendentry{both sides diagonal}
\addplot[color=black!45, dash dot] table[x=tokens, y=muon] {figs/train_delta.dat};
\addlegendentry{Muon}
\draw[gray!60, dashed] (axis cs:0,0) -- (axis cs:2.69,0);
\end{axis}
\end{tikzpicture}
\caption{Paired per-update training-loss difference against the seed-matched
baseline, averaged in 64-update blocks. Because arms consume identical tokens
in identical order, these differences are exact pairings and resolve effects
two orders of magnitude below the sparse evaluations' noise. The orderings
established in the first 200M tokens persist to the end of the budget. Note
that the matched nulls (orange) end below zero despite applying no output
factors at all: that offset is execution path, not geometry, and it is the
quantity Section~\ref{sec:null} subtracts. Their shape differs from the
candidates' --- they start slightly above zero and drift down, while the
candidates open far below it and relax upward.}
\label{fig:traindelta}
\end{figure}

\subsection{A matched null absorbs a fifth to a third of the apparent gain}
\label{sec:null}

The two null arms run their candidate's full machinery --- cotangent capture,
pooled moments, \verb|reduce_scatter|, root computation, host refresh,
scheduling --- and then apply identity factors. If output geometry were the only
thing the candidates changed, the nulls would land on the baseline. They do not.

\nullall{} improves selection loss by $0.001286$ and C4 by $0.002298$ nats over
the baseline while applying no output factors at all. That is \nullsharedense\%
of the dense arm's selection gain and 36\% of its C4 gain. \nullmlp{} improves
C4 by $0.002043$, which is \nullsharemlpcfour\% of the leading arm's C4 gain. The
honest attribution to geometry is therefore candidate minus null:
$-0.004496$ selection and $-0.004081$ C4 for the dense arm, and $-0.005935$
selection and $-0.004035$ C4 for the leading arm.

\begin{finding}[The control is not neutral, and it changes a conclusion]
Running the candidate's code path with identity factors recovers 22--36\% of
the apparent improvement. This is not a small correction to an effect of
0.006 nats. It is large enough to have inverted this project's first
conclusion: the arm favored after the first wave, \texttt{outer\_left\_down},
improved selection loss by $0.001370$, which is \emph{within} the null
level later measured at $0.001286$. Its apparent advantage was substantially
execution path, and nothing in the first wave could have revealed that.
\end{finding}

We do not fully know what the null's gain \emph{is}. The candidate path differs
from the baseline in floating-point association inside the update, in having two
root engines contending for two host worker threads, and in the arrival history
of cached eigenbases. Any of these can perturb a trajectory, and a perturbation
that happens to help is not a result --- but it is also not something a
geometry claim may absorb. One feature of the evidence bounds it: the nulls and
the candidates have different temporal shapes. Both nulls begin the run
\emph{worse} than the baseline ($+0.001717$ and $+0.001220$ nats in the first
536M-token window) and drift to a small advantage by the final window, whereas
the candidates open with a large advantage that decays toward a stable offset
(the dense arm goes from $-0.015940$ to $-0.004975$; see
Table~\ref{tab:windows}). A path artifact that accumulates and a geometry
effect that front-loads are at least distinguishable in shape, which is weak
evidence that they have different origins. And a third null, run in a later
systems wave with upstream arithmetic, independently shows the same end-of-run
effect at $-0.002010$ selection, so this is a property of adding the machinery
rather than a one-run artifact.

One limitation is ours to own: \nullmlp{} ran in a different wave from
\mlpdiag{}, so their difference also spans a change in host scheduling and
source revision. The dense pair is same-wave and is the cleaner measurement.
Closing this is the first item in Section~\ref{sec:program}, and it costs one
run.

\subsection{Freshness beats richness}
\label{sec:freshness}

The width-2048 output moments carry 96.2\% of the output
eigendecomposition work (Section~\ref{sec:method-cost}), so the obvious economy
is to attack them. What we did not expect is that attacking them \emph{improves}
the loss.

\mlpdiag{} replaces exactly those two dense roots with elementwise diagonals and
beats the fully dense arm on selection loss ($-0.006584$ against $-0.005782$)
while costing 17 percentage points less per step. The two arms that retain dense
width-2048 roots, \dense{} and \texttt{outer\_left\_all}, are the two slowest
arms in the study. Meanwhile dropping those factors entirely, in
\texttt{inner\_outer\_hidden}, reaches only $-0.003068$ --- worse than
diagonalizing them.

\begin{finding}[Scale is worth tracking; rotation at this refresh rate is not]
A cheap, exact, always-current approximation of the widest output moments beats
an expensive, stale one, and both beat omitting them. The moments carry real
information, and the part of that information which was not paying for itself
is the eigenbasis, because the eigenbasis cannot be re-estimated fast enough to
keep up with the moment it describes.
\end{finding}

Figure~\ref{fig:freshness} is the direct evidence, taken from the runs' own
counters rather than inferred. The dense-output arms sit at a mean maximum basis
age near three updates against a nominal four-update refresh cadence, and
accumulate roughly 32{,}000 exact-eigendecomposition fallbacks over a run, about
six of 224 statistics per update. The fully diagonal arm logs an age of zero and
\emph{no} fallbacks, because elementwise powers of the current moment need no
basis at all. The dense roots are thus routinely approximations of a moment
that has already moved, while the EMA's own averaging window is only about 16
updates, so a three-update-old basis is a meaningful fraction of the statistic's
memory.

\begin{figure}[t]
\centering
\begin{tikzpicture}
\begin{groupplot}[
  group style={group size=2 by 1, horizontal sep=1.5cm},
  width=0.47\textwidth, height=5.4cm,
  grid=major, grid style={gray!18},
  axis line style={gray!55}, tick align=outside,
  label style={font=\small}, tick label style={font=\footnotesize},
  legend style={font=\scriptsize, draw=gray!40, fill=white,
                fill opacity=0.92, text opacity=1},
  every axis plot/.append style={line width=0.9pt},
]
\nextgroupplot[xlabel={Tokens (billions)}, ylabel={Max output basis age (updates)},
  xmin=0, xmax=2.69, ymin=-0.2, ymax=5,
  legend style={at={(0.98,0.04)}, anchor=south east}]
\addplot[color=laBlueLine!55] table[x=tokens, y=inner_outer_all_age] {figs/freshness.dat};
\addlegendentry{dense output}
\addplot[color=laBlueLine] table[x=tokens, y=inner_outer_mlpdiag_age] {figs/freshness.dat};
\addlegendentry{dense at 512 only}
\addplot[color=laPurpleLine] table[x=tokens, y=inner_outer_leftdiag_age] {figs/freshness.dat};
\addlegendentry{all diagonal}

\nextgroupplot[xlabel={Tokens (billions)}, ylabel={Cumulative exact fallbacks},
  xmin=0, xmax=2.69, ymin=0, ymax=34000,
  scaled y ticks=false, yticklabel style={/pgf/number format/fixed},
  legend style={at={(0.03,0.97)}, anchor=north west}]
\addplot[color=laBlueLine!55] table[x=tokens, y=inner_outer_all_fb] {figs/freshness.dat};
\addlegendentry{dense output}
\addplot[color=laBlueLine] table[x=tokens, y=inner_outer_mlpdiag_fb] {figs/freshness.dat};
\addlegendentry{dense at 512 only}
\addplot[color=laPurpleLine] table[x=tokens, y=inner_outer_leftdiag_fb] {figs/freshness.dat};
\addlegendentry{all diagonal}
\end{groupplot}
\end{tikzpicture}
\caption{Why the cheap approximation wins. Left: the maximum age of an output
eigenbasis in use, averaged over 64-update blocks. Dense-output arms hold bases
about three updates old against a four-update refresh target, while the
diagonal arm recomputes exactly every update and logs age zero. Right:
cumulative exact-eigendecomposition fallbacks, which occur when the cached
basis cannot be repaired within tolerance --- roughly 32{,}000 per run with
dense output factors, 20{,}000 when only the width-512 roots stay dense, and
none at all when every output factor is diagonal.}
\label{fig:freshness}
\end{figure}

Two observations clarify this result, and Section~\ref{sec:mechanism} resolves
the underlying mechanism directly on the saved checkpoints. First, richness and
freshness vary together in Table~\ref{tab:main}, because diagonal factors are
both coordinate-restricted and age-zero; auditing the width-2048 pre-SwiGLU
moments $B_{\text{gate}}$ and $B_{\text{up}}$ in Section~\ref{sec:mech-hadamard}
shows that their off-diagonal correlations already sit at the finite-sample
noise floor (\mechOutGateCorrRmsMed{} and \mechOutUpCorrRmsMed{}), so diagonal
scaling captures \mechOutGateExpVarLPct\% and \mechOutUpExpVarLPct\% of
$\|L - I\|_F^2$ without estimating noisy off-diagonal rotations. Second, a later
wave pinned input-basis installation to a fixed age of one update, verifying
zero discarded refreshes across all four of its arms, and the diagonal-output
arms still beat the baseline by $0.0022$--$0.0038$ nats of final C4 there,
removing variable \emph{input} staleness as a confound.

The latency variants point the same way without resolving it. Lagging the dense
left factors by one update costs $+0.000113$ nats of paired training loss
(s.e.\ $0.000043$) and $+0.001091$ mean late C4 --- small but resolvable, so a
one-update lag is close to free on loss but not free. Lag was introduced to
hide the post-backward wait, and it did not deliver: the lagged dense arm's
end-to-end overhead is $+31.3$\%, worse than the unlagged arm's own cost profile
once other changes landed in the same wave, so we cannot attribute timing
changes to lag alone.

\subsection{Replication across seeds and one batch size}
\label{sec:replication}

The leading geometry was run on three seeds with paired initialization and data
order within each seed. On the fixed external C4 set, the metric that is
comparable across seeds, it improves final loss by $0.006264$, $0.003738$ and
$0.003651$ nats on seeds 0, 1 and 2 --- mean $0.004551$, and negative on every
seed. The 11-point late means are $-0.006078$, $-0.004153$ and $-0.003967$.
Figure~\ref{fig:seeds} shows the spread against the null level.

\begin{figure}[t]
\centering
\begin{tikzpicture}
\begin{axis}[
  width=0.78\textwidth, height=4.6cm,
  xlabel={Seed}, ylabel={$\Delta$ C4 (nats)},
  xtick={0,1,2}, xticklabels={0 (selection seed), 1, 2},
  xmin=-0.5, xmax=2.5, ymin=-0.0075, ymax=0.0012,
  grid=major, grid style={gray!18},
  axis line style={gray!55}, tick align=outside,
  scaled y ticks=false,
  yticklabel style={/pgf/number format/fixed, /pgf/number format/precision=3},
  label style={font=\small}, tick label style={font=\footnotesize},
  legend columns=2,
  legend style={at={(0.99,0.97)}, anchor=north east, font=\scriptsize,
                draw=gray!40, fill=white, column sep=6pt},
]
\addplot[only marks, mark=*, mark size=2.6pt, color=laBlueLine]
  table[x expr=\thisrow{seed}, y=dc4_final] {figs/seeds.dat};
\addlegendentry{final C4}
\addplot[only marks, mark=square*, mark size=2.3pt, color=laPurpleLine]
  table[x expr=\thisrow{seed}, y=dc4_late] {figs/seeds.dat};
\addlegendentry{late-11 mean C4}
\draw[gray!60, dashed] (axis cs:-0.5,0) -- (axis cs:2.5,0);
\draw[laOrangeLine!75, dashed] (axis cs:-0.5,-0.002043) -- (axis cs:2.5,-0.002043);
\node[font=\tiny, text=laOrangeLine!80!black, anchor=west]
  at (axis cs:-0.45,-0.00255) {matched-null C4 level};
\end{axis}
\end{tikzpicture}
\caption{The leading geometry on three seeds, measured on the fixed C4 set.
The effect is negative on every seed but varies by nearly a factor of two, and
seed 0 --- the largest --- is also the seed on which the geometry was selected.
The orange line is the matched null's C4 gain: on seeds 1 and 2 the candidate's
margin over that control is roughly 0.002 nats, which is the honest size of the
replicated geometry effect.}
\label{fig:seeds}
\end{figure}

Three qualifications matter more than the mean. Seed 0 is the seed on which
this geometry was selected among eight candidates, and it is also the seed with
the largest effect; the two fresh seeds give 60\% and 58\% of its margin. The
narrow-window behaviour differs by seed: on seed 1 the advantage shrinks
monotonically as the learning rate decays toward zero ($-0.004749$ over the
late 11 evaluations against $-0.002269$ over the final 4), while on seed 0 it
does not. And no baseline seed-to-seed variability estimate exists for this
recipe, so we cannot convert three paired runs into a confidence statement. We
report the three numbers and decline to put an interval on them.

The single larger-batch pair we ran is consistent and smaller: at a 1M-token
batch over the same 2.68B tokens, the leading geometry improves late C4 by
$0.003522$ and selection by $0.002029$ nats at $+14.5$\% end-to-end time. One
pair at one batch size is not a scaling result, and we draw no trend from two
points.

\subsection{The cost side, and why there is still no win}
\label{sec:cost}

Every arm in Table~\ref{tab:main} that improves loss also costs more time and
more memory. The cheapest candidate that clearly clears its null, \alldiag{},
costs $+10.4$\% median step time, $+8.2$\% end-to-end, and 1.8\,GiB of extra
peak allocation. The leading arm costs $+21.9$\% step time. Only two arms are
cheaper than the baseline in time, \texttt{muon} and \texttt{down\_hybrid}, and
both are worse in loss.

We checked whether the fixed-token gains are large enough to pay for the extra
time by asking when each run first reached a given C4 level. They are not. At
targets of 3.59, 3.52 and 3.46 nats, the baseline crosses at 4721, 5566 and 5848
seconds; the leading arm crosses at 5624, 6309 and 6982. Every candidate we
checked, in every wave, reaches every threshold \emph{later} in wall-clock time
than the baseline, including the arms that reach it in fewer tokens. Two arms
reach the 3.52 evaluation one token-interval earlier than the baseline and
still take longer to get there.

\begin{finding}[No elapsed-time advantage exists at any checked target]
Across all completed waves, no output-side arm reaches any loss target in less
wall-clock time than input-side MuonT. The fixed-token improvements are real
and the time costs are larger, so on a fixed time budget the baseline is ahead.
\end{finding}

These crossings reflect the unoptimized research runner rather than the
intrinsic cost of output normalization. Two arms (\mlpdiag{} and \alldiag{}, as
well as \texttt{inner\_outer\_mlpdiag\_lag}) already cross the $3.52$-nat C4
threshold \sysTokenSavingPct\% earlier in tokens than \muont{} ($2.36\text{B}$
versus $2.49\text{B}$ tokens) and lose in elapsed time only because the legacy
runner adds $+18\%$--$+22\%$ per update. Section~\ref{sec:systems} shows that
eliminating the per-step global synchronization barrier, synchronous fallback
stalls, and per-linear custom-op backward wrappers reduces
\texttt{inner\_outer\_mlpdiag\_lag} to $+\sysFastMlpVsMuonLo\%$--$\sysFastMlpVsMuonHi\%$
over vanilla Muon ($+\sysFastMlpVsMuontLo\%$--$\sysFastMlpVsMuontHi\%$ over
\muont{}), crossing the wall-clock break-even threshold (Appendix~\ref{app:numerics}
separately reports a seven-step Newton--Schulz plateau shared by all arms).
